\documentclass[UTF8,12pt,a4paper]{ctexart}

\usepackage{amsmath}
\usepackage{cases}
\usepackage{cite}
\usepackage{graphicx}
\usepackage{enumerate}
\usepackage{algorithm}
\usepackage{caption}
\usepackage[noend]{algpseudocode}
\makeatletter
\renewcommand{\fnum@algorithm}{\fname@algorithm}
\makeatother
\renewcommand{\algorithmicrequire}{\textbf{Input:}}
\renewcommand{\algorithmicensure}{\textbf{Output:}}
\usepackage[margin=1in]{geometry}
\geometry{a4paper}
\setlength{\headheight}{14.5pt}
\usepackage{fancyhdr}
\pagestyle{fancy}
\fancyhf{}
\usepackage{xcolor}
\usepackage{listings}
\lstset{
	basicstyle=\ttfamily\small,
	breaklines,
	tabsize=4,
	extendedchars=false,
	columns=fixed,
	numbers=left,
	frame=single,
	backgroundcolor=\color[RGB]{245,245,244},
	keywordstyle=\color[RGB]{40,40,255},
	numberstyle=\footnotesize\color{darkgray},
	commentstyle=\it\color[RGB]{0,96,96},
	stringstyle=\rmfamily\slshape\color[RGB]{128,0,0},
	showstringspaces=false,
	xleftmargin=2em,
	language=c++,
}

\title{算法原理 Lab3 模拟退火算法}
\author{
	姓名：\underline{冯峻}~~~~~~
	学号：\underline{523031910148}~~~~~~}
\date{\today}
\pagenumbering{arabic}

\begin{document}

\fancyhead[L]{冯峻}
\fancyhead[C]{模拟退火算法}
\fancyfoot[C]{\thepage}

\maketitle
\tableofcontents
\newpage
\section{摘要}

本实验利用模拟退火（Simulated Annealing, SA）算法解决标准单元版图中晶体管布局优化问题。实验核心任务是在已提供的算法框架基础上，补全模拟退火主循环的关键代码，包括：随机扰动生成、合法性检测调用、目标函数评分、Metropolis接受准则实现以及降温策略。

实验完成了\texttt{sa\_core.cpp}中所有TODO部分的代码填空，并额外实现了退火过程数据记录与可视化功能（fj\_pro扩展）。通过多组参数实验，验证了模拟退火算法在标准单元布局问题上的有效性，并分析了起始温度、降温因子、迭代次数等参数对优化结果的影响。

\newpage
\section{实验步骤}

\subsection{模拟退火算法原理}

模拟退火算法是一种启发式全局优化方法，其灵感来源于固体退火过程。算法通过引入随机性来避免陷入局部最优解，其核心思想如下：

\begin{enumerate}
	\item \textbf{初始化}：设定初始温度$T_0$、终止温度$T_t$、降温因子$\alpha$
	\item \textbf{扰动生成}：在当前解的邻域内随机产生一个新解
	\item \textbf{接受准则}：根据Metropolis准则决定是否接受新解
	\item \textbf{降温}：按照降温策略更新温度
	\item \textbf{终止}：当温度降至终止温度时，算法结束
\end{enumerate}

Metropolis接受准则是模拟退火的核心，其数学表达式为：
\[
P(\text{接受新解}) = 
\begin{cases}
1 & \text{if } \Delta E > 0 \text{ （新解更优）} \\
e^{\Delta E / T} & \text{if } \Delta E \leq 0 \text{ （新解较差）}
\end{cases}
\]
其中$\Delta E = E_{\text{new}} - E_{\text{old}}$是新旧解的能量差（本实验中为得分差，得分越高越好，因此当$\Delta E > 0$时直接接受，当$\Delta E < 0$时，根据当前温度情况概率接受），$T$为当前温度。

\subsection{模拟退火算法伪代码}

\begin{algorithm}[H]
	\caption{模拟退火算法 Simulated Annealing}
	\begin{algorithmic}[1]
		\Require 初始解$S_0$，初始温度$T_0$，终止温度$T_t$，降温因子$\alpha$，每温度迭代次数$times$
		\Ensure 优化后的解$S_{best}$
		\State $S_{\text{cur}} \gets S_0$, $T \gets T_0$
		\While{$T > T_t$}
			\For{$i = 1$ to times}
				\State $S_{\text{new}} \gets$ 在$S_{\text{cur}}$邻域内随机生成新解
				\If{$S_{\text{new}}$合法}
					\State $\Delta E \gets \text{Score}(S_{\text{new}}) - \text{Score}(S_{\text{cur}})$
					\If{$\Delta E > 0$ \textbf{or} $\text{random}(0,1) < e^{\Delta E/T}$}
						\State $S_{\text{cur}} \gets S_{\text{new}}$ \Comment{接受新解}
					\EndIf
				\EndIf
			\EndFor
			\State $T \gets \alpha T$ \Comment{降温}
		\EndWhile
		\State \Return $S_{cur}$
	\end{algorithmic}
\end{algorithm}

\subsection{代码补全说明}

本实验的主要内容是在\texttt{sa\_core.cpp}文件中补全模拟退火主函数框架\texttt{simulated\_annealing\_cpp}的关键部分。

\subsubsection{TODO 1: 随机选择位置并调用is\_legal}

从占用位置列表\texttt{occupy}和空位置列表\texttt{empty}中各随机选择一个位置，调用\texttt{is\_legal}函数判断将MOS管从占用位置（即occupy[i]）移动到空位置（即empty[j]）是否合法。is_legal()函数内部会自动尝试各种翻转组合来找到合法布局。
\begin{lstlisting}
// 创建随机分布
std::uniform_int_distribution<int> dist_occ(0, 
    static_cast<int>(occupy.size()) - 1);
std::uniform_int_distribution<int> dist_emp(0, 
    static_cast<int>(empty.size()) - 1);
int i = dist_occ(rng);  // 随机选择占用位置索引
int j = dist_emp(rng);  // 随机选择空位置索引

bool legal;
vector<Mos> pary, nary;
vector<optional<Mos>> pplace, nplace;

if (tp) {  // PMOS行
    std::tie(legal, pary, nary, pplace, nplace) = 
        is_legal(cur_pmos, cur_nmos, cur_pp, cur_np, 
                 occupy[i], empty[j]);
	// 对于 PMOS 行，传入顺序是 (pmos, nmos, pp, np, occupy[i], empty[j])
} else {   // NMOS行，需要交换参数顺序
    std::tie(legal, nary, pary, nplace, pplace) = 
        is_legal(cur_nmos, cur_pmos, cur_np, cur_pp, 
                 occupy[i], empty[j]);
	// 对于 NMOS 行，传入顺序是 (nmos, pmos, np, pp, occupy[i], empty[j])
}
\end{lstlisting}

\subsubsection{TODO 2: 计算当前解和候选解的分数}

使用\texttt{evaluator}函数分别计算当前解和候选解的得分。所有参数按照形式参数列表传入即可。

在计算新解的分数时，要使用更新后的MOS数组\texttt{new\_pary}、\texttt{new\_nary}，以及布局数组\texttt{new\_pp2}和\texttt{new\_np2}。

\begin{lstlisting}
// 计算当前解的分数
auto eval0 = evaluator(w_ref, nets, pins, 
    cur_pmos, cur_nmos, cur_pp, cur_np);
double s0 = std::get<0>(eval0);  // 总分

// 计算候选解的分数
auto eval1 = evaluator(w_ref, nets, pins, 
    new_pary, new_nary, new_pp2, new_np2);
double s1 = std::get<0>(eval1);
\end{lstlisting}

\subsubsection{TODO 3: Metropolis接受准则}

实现Metropolis接受准则，决定是否接受候选解。这里dist01(rng)用于随机生成一个0到1之间的浮点数，用作概率与当前温度下的阈值概率比较。

\begin{lstlisting}
double delta = s1 - s0;  // 得分差，正值表示新解更好
bool accept = false;

if (delta > 0) {
    // 新解更好，直接接受
    accept = true;
} else {
    // 新解更差，以概率e^(delta/t)接受
    double threshold = std::exp(delta / t);
    if (dist01(rng) < threshold) {
        accept = true;
    }
}
\end{lstlisting}

\subsubsection{TODO 4: 接受后更新当前解}

当新解被接受时，更新当前解并在行末尾追加空位。

\begin{lstlisting}
if (accept) {
    cur_pmos = new_pary;
    cur_nmos = new_nary;
    cur_pp = new_pp2;
    cur_np = new_np2;
    
    // 追加空位，为后续扰动预留空间
    cur_pp.push_back(std::nullopt);
    cur_pp.push_back(std::nullopt);
    cur_np.push_back(std::nullopt);
    cur_np.push_back(std::nullopt);
    
    // 行切换逻辑
    if (tp) {
        if (!cur_nmos.empty()) tp = !tp;
    } else {
        if (!cur_pmos.empty()) tp = !tp;
    }
}
\end{lstlisting}

\subsubsection{TODO 5: 降温}

采用指数降温策略：

\begin{lstlisting}
t = t * decrease;  // 指数降温
\end{lstlisting}

\newpage
\section{实验扩展工作}

除了完成必要的TODO部分，还额外实现了以下扩展功能：

\subsection{退火过程数据记录}

在\texttt{simulated\_annealing\_cpp}函数中添加了数据记录功能，将每次迭代的温度、当前解分数、候选解分数以及是否接受（0/1标识）记录到CSV文件。

要完成这一步，有如下代码片段需要改动：

\textbf{1. 在\texttt{simulated\_annealing\_cpp}函数参数列表中添加一个字符串参数\texttt{data\_filename}，表示数据文件名。}
\begin{lstlisting}
tuple<vector<Mos>, vector<Mos>, vector<optional<Mos>>, vector<optional<Mos>>>
simulated_annealing_cpp(
	int w_ref,  // fj: 参考宽度
	const Nets& nets,  // fj: net 名称数组
	const Pins& pins,  // fj: pin 名称数组
	const vector<Mos>& pmos_ary, // fj: 初始 PMOS 数组
	const vector<Mos>& nmos_ary, // fj: 初始 NMOS 数组
	const vector<optional<Mos>>& pp_ary, // fj: 初始 PMOS 布局行
	const vector<optional<Mos>>& np_ary, // fj: 初始 NMOS 布局行
	double t0,        // 起始温度
	double tt,        // 终止温度（停止条件：t <= tt）
	double decrease,  // 降温因子，例如 0.95，每轮 t *= decrease
	int times,         // 每个温度下尝试的 move 次数
	std::string data_filename // fj: 记录退火过程中得分的数据文件名字
) 
\end{lstlisting}

\textbf{2. 退火过程中，每一次move操作都记录当前解分数以及候选解分数。}
\begin{lstlisting}
// fj_pro: 初始化数据文件，用于记录退火过程中的分数变化
std::ofstream data_file;
bool logging = !data_filename.empty();
if (logging) {
	data_file.open(data_filename);
	// 写入表头: 温度, 当前分数, 候选分数, 是否接受(1/0)
	data_file << "Temp,CurrentScore,CandidateScore,Accepted\n";
}

// fj_pro: 计算初始分数 (假设你代码里有这部分，如果没有，需要在循环前算一次)
auto eval_init = evaluator(w_ref, nets, pins, cur_pmos, cur_nmos, cur_pp, cur_np);
double s_init = std::get<0>(eval_init);  // 初始解的总分

std::uniform_real_distribution<double> dist01(0.0, 1.0);

// ......

// fj_pro: 将当前温度、当前分数、候选分数、是否接受记录到数据文件中，用于绘图
if (logging) {
	data_file << t << "," << s0 << "," << s1 << "," << (accept ? 1 : 0) << "\n";
}
\end{lstlisting}

\textbf{3. 在\texttt{sa\_bindings.cpp}中，将\texttt{simulated\_annealing\_cpp}函数绑定到Python接口时，添加新的字符串参数。}
\begin{lstlisting}
// fj_pro: 修改为带有 data_filename 参数的 run_sa
m.def("run_sa", &run_sa, "Run simulated annealing core",
	py::arg("w_ref"),
	py::arg("nets"),
	py::arg("pins"),
	py::arg("pmos_ary"),
	py::arg("nmos_ary"),
	py::arg("pp_ary"),
	py::arg("np_ary"),
	py::arg("t0"),
	py::arg("tt"),
	py::arg("decrease"),
	py::arg("times"),
	py::arg("data_filename") = "" // [新增] 允许 Python 调用时不传此参数
);
\end{lstlisting}

\textbf{4. 在eda.py中调用模拟退火算法时，需要传入数据文件名参数。}
\begin{lstlisting}[language=Python]
# fj_pro: 将sys.argv修改为更方便的argparse处理，同时加入数据文件
parser = argparse.ArgumentParser(description="Simulated Annealing Placement Tool")
parser.add_argument("netlist", help="Path of the netlist file")
parser.add_argument("cell_name", help="Name of the cell to process")
parser.add_argument("data_file", help="File to log annealing data")
args = parser.parse_args()
cells = args.netlist
cell = args.cell_name
data_file = args.data_file
\end{lstlisting}

\subsection{迭代曲线绘制脚本}

编写了Python脚本\texttt{plot\_iter.py}，用于可视化退火过程中的收敛曲线。该脚本读取CSV数据文件，绘制历史最佳分数随温度变化的曲线。同时，该脚本接受\texttt{x_mode}参数，支持绘图时更改横坐标为对数模式（指数速率降温）或者线性模式（线性速率降温）。

\begin{lstlisting}[language=Python]
# 计算历史最佳分数曲线，该曲线必定连续，能够很直白的看出算法优化效果
df['BestScore'] = df['CurrentScore'].cummax()

plt.figure(figsize=(10, 6))

# 1. 绘制当前分数 (使用散点图，更能体现高温下的随机性)
# plt.scatter(df['Temp'], df['CurrentScore'], s=10, c='blue', alpha=0.15, label='Current Score')
# 如果绘制处所有数据点，则svg文件也会很大，所以直接放弃散点图

# 2. 绘制历史最佳分数 (使用连线，体现收敛趋势)
plt.plot(df['Temp'], df['BestScore'], color='red', linewidth=2, label='Best Score')

# ......

# 设置 X 轴为对数坐标 (因为退火温度通常是指数下降的，线性坐标会导致低温段挤在一起)
if x_mode == 'log':
	plt.xscale('log')
\end{lstlisting}

\subsection{其他}

由于后续会将程序在终端打印出的信息保存到文件，为了防止中文乱码，在\texttt{set_up.py}中，做出如下更改
\begin{lstlisting}[language=Python]
extra_compile_args = []
if sys.platform == "win32":
	extra_compile_args = ["/std:c++17", "/O2", "/utf-8"]
\end{lstlisting}

增加了utf-8参数，否则在执行命令\texttt{python setup.py build_ext --inplace}时会出现：“warning C4819: 该文件包含不能在当前代码页(936)中表示的字符”。因为源文件是 UTF-8 编码，但 Windows 的中文命令行默认使用 GBK (代码页 936)。

\newpage
\section{实验过程及数据}

\subsection{实验环境}

\begin{itemize}
	\item \textbf{操作系统}：Windows 11 
	\item \textbf{编程语言}：C++17、Python 311.13
	\item \textbf{编译工具}：pybind11用于C++与Python绑定
	\item \textbf{测试单元}：AN2D2（8个MOS管：4个PMOS + 4个NMOS）
\end{itemize}

\subsection{编译与运行流程}
在加入上述扩展工作后，编译与运行流程发生如下变化：

\textbf{1. 编译}

\texttt{python setup.py build\_ext --inplace}：编译C++扩展模块。

该命令会告诉 setup.py 把编译好的 .pyd 文件直接放在当前源代码目录下，这样 eda.py 就能直接引用。

\textbf{2. 运行模拟退火算法}

\texttt{python eda.py cells.spi AN2D2 "results\20000_1_88_2000\data.csv" > "results\20000_1_88_2000\process.log"}

将eda.py打印的退火日志信息保存到log文件，将退火过程中的评分保存到.csv文件， data_file用于记录模拟退火过程中的数据文件。

\textbf{3. 评分}

\texttt{python evaluator.py AN2D2.json AN2D2 cells.spi > "results\20000_1_88_2000\score.log"}

将评分的信息保存到.log文件。

\textbf{4. 绘制迭代曲线}

\texttt{python plot\_iter.py "results\20000_1_88_2000\data.csv" log}

\subsection{实验参数设置}

实验测试了多组参数组合，参数命名规则为\texttt{x\_y\_z\_w}，其中x, y, z, w的含义分别是，在当前文件夹保存的退火模式中：
\begin{enumerate}
	\item 起始温度$T_0$： $x \times \text{mos\_num}$，其中mos\_num为当前单元的MOS管数量
	\item 终止温度$T_t$：$0.y$，在本次实验中固定为0.1
	\item 降温因子$\alpha$：$0.z$，如88表示0.88
	\item 每温度迭代次数$times$：$w \times \text{mos\_num}$
\end{enumerate}

\subsection{实验数据}

以AN2D2标准单元（mos\_num=8）为例，各参数组合的实验结果如下：

\begin{center}
\begin{tabular}{|c|c|c|c|c|}
	\hline
	\textbf{参数组合} & \textbf{$T_0$} & \textbf{$\alpha$} & \textbf{$times$}  & \textbf{备注} \\
	\hline
	20000\_1\_88\_2000 & 160000 & 0.88 & 16000 & 高温慢降，长迭代 \\
	\hline
	20000\_1\_88\_1000 & 160000 & 0.88 & 8000 & 高温慢降，短迭代 \\
	\hline
	20000\_1\_7\_2000 & 160000 & 0.70 & 16000 & 高温快降，长迭代 \\
	\hline
	20000\_1\_7\_1000 & 160000 & 0.70 & 8000 & 高温快降，短迭代 \\
	\hline
	20000\_1\_5\_2000 & 160000 & 0.70 & 16000 & 高温快降，长迭代 \\
	\hline
	20000\_1\_5\_1000 & 160000 & 0.70 & 8000 & 高温快降，短迭代 \\
	\hline
	10000\_1\_88\_2000 & 160000 & 0.70 & 8000 & 中温快降 \\
	\hline
	10000\_1\_8\_1000 & 80000 & 0.80 & 8000 & 低温中速降 \\
	\hline
\end{tabular}
\end{center}

\subsection{收敛曲线示例}

通过fj\_pro扩展功能生成的退火收敛曲线保存在\texttt{results/}目录下的各参数文件夹中。曲线展示了历史最佳分数随温度下降的变化趋势，体现了模拟退火算法的收敛特性。

例如，在退火模式：\texttt{results/20000\_1\_88\_2000/}下，生成的收敛曲线如下所示：

\begin{center}
	\centering
	\includegraphics[width=0.8\textwidth]{results/20000_1_88_2000/data.pdf}
	\caption{参数组合20000\_1\_88\_2000的收敛曲线}
	\label{fig:convergence_curve}
\end{center}

\section{分析与讨论}

\subsection{参数影响分析}

\subsubsection{降温因子$\alpha$的影响}

\begin{itemize}
	\item \textbf{较大的$\alpha$（如0.88）}：降温较慢，算法有更多时间在每个温度水平进行探索，有利于找到更优解，但运行时间较长。
	\item \textbf{较小的$\alpha$（如0.70）}：降温较快，算法更快收敛，但可能过早陷入局部最优。
\end{itemize}

\subsubsection{每温度迭代次数$times$的影响}

\begin{itemize}
	\item \textbf{较大的$times$}：每个温度下进行更多扰动尝试，增加找到更优解的概率，但增加计算时间。
	\item \textbf{较小的$times$}：减少计算时间，但可能导致搜索不充分。
\end{itemize}

\subsubsection{起始温度$T_0$的影响}

\begin{itemize}
	\item \textbf{较高的$T_0$}：初期接受差解的概率高，有利于跳出局部最优，探索范围更广。
	\item \textbf{较低的$T_0$}：初期即表现出较强的贪心性质，可能错过全局最优。
\end{itemize}

\subsection{SA与贪心算法的效果差异}

\textbf{【待完成】}

此部分需要实现简单贪心算法（只接受改进解）进行对比实验，分析以下内容：

\begin{enumerate}
	\item 贪心算法的实现方式（Metropolis准则中只接受$\Delta E > 0$的情况）
	\item 在相同测试用例上，SA与贪心的最终得分对比
	\item SA相对于贪心在避免局部最优方面的优势分析
	\item 运行时间对比
\end{enumerate}

% 预留空间用于后续补充实验数据
\begin{center}
\begin{tabular}{|c|c|c|c|}
	\hline
	\textbf{算法} & \textbf{最终得分} & \textbf{运行时间} & \textbf{备注} \\
	\hline
	模拟退火 & --- & --- & 待补充 \\
	\hline
	贪心算法 & --- & --- & 待补充 \\
	\hline
\end{tabular}
\end{center}

\subsection{实验总结与收获}

\begin{enumerate}
	\item \textbf{理解了模拟退火算法的核心思想}：通过引入随机性（Metropolis准则）来避免陷入局部最优，在高温时广泛探索，低温时精细收敛。
	
	\item \textbf{掌握了模拟退火的完整流程}：邻域产生$\rightarrow$合法性检查$\rightarrow$目标函数评分$\rightarrow$接受准则$\rightarrow$降温策略。
	
	\item \textbf{理解了参数调节的重要性}：不同的参数组合会显著影响算法的收敛速度和最终解的质量，需要根据具体问题进行调优。
	
	\item \textbf{实践了C++与Python的混合编程}：通过pybind11将计算密集型的C++代码封装为Python模块，兼顾了性能和易用性。
\end{enumerate}

\end{document}
